% ==============================================================================
% CONFIGURAZIONI E PACCHETTI DELLA TESI
% ==============================================================================

% Layout e margini pagina (Linee Guida DISI UniTn)
\usepackage[paperheight=29.7cm,paperwidth=21cm,top=2cm,bottom=2cm,outer=1.5cm,inner=2.5cm]{geometry}
\usepackage{setspace}
\singlespacing
\setlength{\parindent}{0pt}
\setlength{\parskip}{3pt plus 1pt minus 1pt}
\raggedbottom

% Gestione elenchi (senza rientro a sinistra, allineati al margine del testo)
\usepackage[inline]{enumitem}
\setlist{leftmargin=*, topsep=4pt, itemsep=2pt, partopsep=0pt, parsep=0pt}

% Titoli capitoli e sezioni in linea (stile tesi esempio)
\usepackage{titlesec}
\titleformat{\chapter}[hang]
  {\normalfont\huge\bfseries}{\thechapter}{1em}{\huge}
\titlespacing*{\chapter}{0pt}{15pt}{15pt}
\titlespacing*{\section}{0pt}{14pt}{6pt}
\titlespacing*{\subsection}{0pt}{10pt}{4pt}

\usepackage{plain}
\usepackage{graphicx}
\usepackage{subcaption}
\usepackage{float}
\usepackage{booktabs}

% Font e tipografia
\usepackage[T1]{fontenc}
\usepackage{newtxtext,newtxmath} % Font Times New Roman (versione moderna per LaTeX)

% Supporto lingua italiana e codifica
\usepackage[utf8x]{inputenc}
\usepackage[italian]{babel}

% Colori e grafica TikZ
\usepackage{xcolor}
\usepackage{standalone}
\usepackage{tikz}
\usetikzlibrary{
    shapes.geometric,
    arrows.meta,
    positioning,
    fit,
    backgrounds,
    calc,
    shadows.blur,
    decorations.pathreplacing
}

% Palette multicolore convenzionale accademica (Standard Scientific / ColorBrewer)
\definecolor{skblue_head}{RGB}{31, 119, 180}     % Blu classico (Dataset)
\definecolor{skblue_bg}{RGB}{224, 238, 250}
\definecolor{skblue_txt}{RGB}{12, 45, 82}

\definecolor{skteal_head}{RGB}{23, 145, 155}     % Ottanio / Teal (Feature & Encoding)
\definecolor{skteal_bg}{RGB}{220, 244, 246}
\definecolor{skteal_txt}{RGB}{10, 52, 58}

\definecolor{skgreen_head}{RGB}{44, 155, 45}     % Verde classico (Modelli & Addestramento)
\definecolor{skgreen_bg}{RGB}{225, 245, 225}
\definecolor{skgreen_txt}{RGB}{15, 58, 20}

\definecolor{skorange_head}{RGB}{225, 110, 20}    % Arancione / Ambra (Valutazione)
\definecolor{skorange_bg}{RGB}{254, 235, 220}
\definecolor{skorange_txt}{RGB}{85, 35, 10}

\definecolor{skpurple_head}{RGB}{135, 85, 175}   % Viola classico (Spiegabilità)
\definecolor{skpurple_bg}{RGB}{242, 230, 248}
\definecolor{skpurple_txt}{RGB}{50, 18, 70}

% ==============================================================================
% COMANDI PERSONALIZZATI (CUSTOM MACROS)
% ==============================================================================

% Inclusione figure TikZ standalone scalate
% Sintassi: \inputtikz[larghezza (default \textwidth)]{file.tex}
\newcommand{\inputtikz}[2][\textwidth]{\resizebox{#1}{!}{\input{#2}}}

% ==============================================================================
% COLLEGAMENTI IPERTESTUALI (HYPERREF)
% ==============================================================================
\usepackage[
    hidelinks,                  % Rende i link cliccabili senza riquadri o colori invasivi
    linktoc=all,                % Rende sia il testo che i numeri di pagina cliccabili nell'indice
    unicode=true,               % Supporto caratteri unicode nei segnalibri PDF
    pdfstartview={FitH},        % Adatta alla larghezza della pagina all'apertura
    bookmarks=true,             % Abilita i segnalibri PDF
    bookmarksnumbered=true      % Numera le sezioni nei segnalibri PDF
]{hyperref}

